\slajdid{08-s017}
\begin{frame}[label=08-s017]{„Dobra“ karakteristika}
\gusttekst\setlength{\parskip}{2pt}
\begin{columns}[c,onlytextwidth]\column{.55\textwidth}\centering\input{slike/tikz/s017_dozvoljene_oblasti.tex}
\column{.43\textwidth}\centering\input{slike/tikz/s017_tri_oblasti.tex}\end{columns}\medskip
Znači dobra karakteristika prenosa ima tri oblasti.\par
1. Oblast logičke nule sa malim pojačanjem, manjim od jedan\par
2. Prelaznu zonu koja treba da je što uža i treba da obezbedi prelaz sa logičke nule na logičku jedinicu. Kako ova zona treba da bude što uža (robusnost, margine šuma) očigledno je da pojačanje u ovoj zoni treba da je što veće, puno veće od jedan.\par
3. Oblast logičke jedinice sa malim pojačanjem, manjim od jedan.
\end{frame}
